\chapter{Ejercicios sobre Clases de Complejidad}

\begin{ejercicio}~
    \begin{enumerate}
        \item[(a)] Un grafo dirigido se dice que es \textit{acíclico} si no tiene ciclos. Demostrar que todo grafo finito dirigido acíclico tiene una \textit{fuente} (un nodo al que no llegan arcos).

        \item[(b)] Utilizar (a) para demostrar que un grafo dirigido con $n$ nodos es acíclico si y solo si se pueden numerar los nodos del $1$ al $n$ de manera que siempre los arcos van desde números más pequeños a números más grandes.

        \item[(c)] Utilizar (a) y (b) para describir un algoritmo en tiempo polinómico para determinar cuándo un grafo dirigido es acíclico.
    \end{enumerate}
\end{ejercicio}

\begin{ejercicio}~
    \begin{enumerate}
        \item[(a)] Demostrar que un grafo es \textit{bipartito} (se puede dividir en dos partes, que no son necesariamente iguales, de manera que todas las aristas van de una parte a otra) si y solo si todos sus ciclos son de longitud par.

        \item[(b)] Utilizar (a) para describir un algoritmo en tiempo polinómico para comprobar si un grafo es bipartito.
    \end{enumerate}
\end{ejercicio}

\begin{ejercicio}
    Demostrar que \Pol \text{ es cerrada para la unión y la intersección}.
\end{ejercicio}

\begin{ejercicio}
    Sea $C$ una clase de funciones de los enteros no negativos en los enteros no negativos. Se dice que $C$ es cerrada para la composición polinómica por la izquierda si $f(n) \in C$ implica que $p(f(n)) = O(g(n))$ para alguna función $g(n) \in C$, donde $p(n)$ es un polinomio cualquiera.
    Se dice que $C$ es cerrada para la composición polinómica por la derecha si $f(n) \in C$ implica que $f(p(n)) = O(g(n))$ para alguna función $g(n) \in C$, donde $p(n)$ es un polinomio cualquiera. \\

    \noindent Determinar cuáles de las siguientes clases de funciones son cerradas para la composición polinómica por la derecha y cuáles lo son por la izquierda.

    \begin{enumerate}
        \item[(a)] $\{n^k : k > 0\}$.
        \item[(b)] $\{n \cdot k : k > 0\}$.
        \item[(c)] $\{k^n : k > 0\}$.
        \item[(d)] $\{2^{n^k} : k > 0\}$.
        \item[(e)] $\{\log^k(n) : k > 0\}$.
        \item[(f)] $\{\log(n)\}$.
    \end{enumerate}

    Nota: en los dos últimos casos se entiende que las funciones se redondean al entero más cercano. En vez de considerar todos los polinomios, como estamos hablando de órdenes de complejidad, será suficiente considerar los polinomios de la forma $p(n) = n^l$, donde $l$ es un número natural.
\end{ejercicio}

\begin{ejercicio}
    Sea $L = \{wcw : w \in \{0,1\}^*\}.$ Demostrar que $L$ se puede reconocer en espacio $\log(n)$.
\end{ejercicio}

\begin{ejercicio}
    Demostrar que si $L$ está en \Pol, entonces $L^*$ también está en \Pol.
\end{ejercicio}

\begin{ejercicio}
    Demostrar que si $L$ está en \NP, entonces $L^*$ también está en \NP.
\end{ejercicio}

\begin{ejercicio}
    Demostrar que \NP \text{ es cerrada para la unión y la intersección}.
\end{ejercicio}

\begin{ejercicio}
    Demostrar que si $\NP \neq \CoNP$, entonces $\Pol \neq \NP$.
\end{ejercicio}

\begin{ejercicio}
    Supongamos que una entrada es una palabra de paréntesis. Demostrar que determinar si están emparejados y anidados correctamente está en \Loga. Lo están $(()())$ y $((()))$, pero no lo están $)()($ ni $(()))($.
\end{ejercicio}

\begin{ejercicio}
    Consideremos el problema anterior, pero ahora con dos tipos de paréntesis: $(,)$ y $[,]$. Demostrar que determinar si están bien escritos 
    también está en \Loga. Aquí $([]())$ está permitido, pero $([)]$ no. Demostrar que si la entrada se lee de izquierda a derecha sin volver hacia atrás, 
    requieren $O(n)$ de memoria.
\end{ejercicio}

\begin{ejercicio}
    Demostrar que el problema del palíndromo requiere $O(n)$ de memoria si no se puede volver hacia atrás.
\end{ejercicio}

\begin{ejercicio}
    Demostrar que $\NP \neq \ESP{n}$. Nota: Este ejercicio es muy complicado. Como sugerencia, para probar que 
    las clases son diferentes, probar que $\NP$ cumple la propiedad de ser cerrada para reducciones en $L$ 
    (si $P_1 \propto P_2$ y $P_2$ está en $\NP$, entonces $P_1$ está en $\NP$), mientras que $\ESP{n}$, no.
\end{ejercicio}

\newpage
\section{Soluciones}
\setcounter{ejercicio}{0}

\begin{ejercicio}~
    \begin{enumerate}
        \item[(a)] Un grafo dirigido se dice que es \textit{acíclico} si no tiene ciclos. Demostrar que todo grafo finito dirigido acíclico tiene 
        una \textit{fuente} (un nodo al que no llegan arcos). \\

            Sea $G=(V,E)$ un grafo dirigido finito acíclico. Por ser $G$ finito, podemos suponer que $|V| = n \in \N$. Supongamos ahora, 
            por reducción al absurdo, que $G$ no tiene ninguna fuente. Entonces, todo nodo de $G$ tiene al menos un arco incidente. Tomamos un nodo cualquiera $v_1 \in V$. Como $v_1$ no es fuente, 
            existe un nodo $v_2 \in V$ tal que $(v_2,v_1)\in E$. De nuevo, como $v_2$ no es fuente, existe $v_3 \in V$ tal que $(v_3,v_2)\in E$. 
            Repitiendo el argumento, obtenemos una sucesión de nodos $$v_1,v_2,\dots,v_{n+1}$$
            tales que, para todo $i=1,\dots,n$,
            $$
            (v_{i+1},v_i)\in E.
            $$

            Como estos $n+1$ nodos pertenecen al conjunto de vértices $V$, que tiene solamente $n$ elementos, existen $1\leq i<j\leq n+1$ tales que 
            $v_i=v_j$. Entonces tenemos el ciclo dirigido
            $$
            v_j \to v_{j-1}\to \cdots \to v_i=v_j,
            $$
            lo cual contradice que $G$ sea acíclico. Por tanto, $G$ tiene al menos una fuente.

        \item[(b)] Utilizar (a) para demostrar que un grafo dirigido con $n$ nodos es acíclico si y solo si se pueden numerar los nodos del $1$ al $n$ 
        de manera que siempre los arcos van desde números más pequeños a números más grandes. \\

            Queremos demostrar que, llamando $P$ a la propiedad de numeración del enunciado, si $G$ es un grado dirigido con $n$ nodos, entonces $$G \text{ es acíclico} \iff \text{se cumple P}$$
            
            \begin{itemize}
                \item[$\Longrightarrow)$] Supongamos primero que $G$ es acíclico. Como $G$ es finito y tiene $n$ nodos, por el apartado (a), $G$ tiene una fuente. Numeramos dicha fuente con el número $1$ y la eliminamos del grafo junto con todos sus arcos salientes.
                    El grafo restante sigue siendo finito y acíclico, por ser este subgrafo de un grafo finito y acíclico. Por tanto, si queda algún nodo, podemos volver a aplicar el apartado (a) y obtener una nueva fuente. La numeramos con el número $2$ y la eliminamos junto con sus arcos salientes.
                    Repitiendo este proceso, como inicialmente había $n$ nodos, tras $n$ pasos todos los nodos quedan numerados con los números $1,\dots,n$. Veamos que esta numeración cumple la propiedad pedida. 
                    Sea $(u,v)\in E$ un arco cualquiera del grafo original. Supongamos que $v$ se numeró antes que $u$. Entonces, en el momento en que se numeró $v$, el nodo $u$ todavía no había sido eliminado. Además, el arco $(u,v)$ tampoco había sido eliminado, pues solo 
                    se eliminan los arcos salientes de los nodos que se van numerando, y $u$ todavía seguía en el grafo. Por tanto, en ese momento $v$ tenía un arco entrante, luego no podía ser fuente, pero si se había eliminado $v$ era porque era fuente, lo cual es una contradicción. 
                    Así, necesariamente $u$ se numera antes que $v$, es decir, $\operatorname{num}(u)<\operatorname{num}(v)$ y se cumple $P$.
                \item[$(\Longleftarrow$] Recíprocamente, supongamos que se cumple $P$, es decir, todos los nodos de $G$ se pueden numerar del $1$ al $n$ de forma que todo arco va desde un nodo con número menor a un nodo con número mayor. Veamos que $G$ es acíclico.
                    Si $G$ tuviera un ciclo dirigido
                    $$
                    v_1\to v_2\to \cdots \to v_k\to v_1,
                    $$
                    para cierto $1 \leq k \leq n$ entonces, como por hipótesis se cumple $P$, tendríamos
                    $$
                    \operatorname{num}(v_1)<\operatorname{num}(v_2)<\cdots<\operatorname{num}(v_k)<\operatorname{num}(v_1) \Longrightarrow \operatorname{num}(v_1) < \operatorname{num}(v_1) 
                    $$
                    lo cual contradice que se cumpla $P$. Por tanto, $G$ no tiene ciclos dirigidos, es decir, $G$ es acíclico.
            \end{itemize}

        \item[(c)] Utilizar (a) y (b) para describir un algoritmo en tiempo polinómico para determinar cuándo un grafo dirigido es acíclico. \\
            
            El algoritmo es el siguiente:

            \begin{algorithm}
                \caption{Decidir si un grafo dirigido es acíclico}
                \label{algorithm:aciclico}
                \begin{algorithmic}
                    \State \textbf{Entrada:} Un grafo dirigido finito $G=(V,E)$

                    \While{$V\neq \emptyset$}
                        \State Buscar un nodo $v\in V$ fuente.
                        \If{no existe tal nodo}
                            \State Rechazar, es decir, responder que $G$ no es acíclico.
                        \EndIf
                        \State Eliminar $v$ del grafo, junto con todos sus arcos salientes.
                    \EndWhile

                    \State Aceptar, es decir, responder que $G$ es acíclico.
                \end{algorithmic}
            \end{algorithm}

            Veamos que el algoritmo es correcto. Si el algoritmo acepta, entonces ha eliminado todos los nodos del grafo. Por tanto, podemos numerar los nodos en el mismo orden en que han sido eliminados. Como en cada paso se elimina un nodo fuente, por el razonamiento del apartado (b), 
            todo arco del grafo original va desde un nodo eliminado antes hacia un nodo eliminado después. Así, se obtiene una numeración de los nodos del $1$ al $n$ en la que todos los arcos van desde números más pequeños a números más grandes. Nuevamente por el apartado (b), $G$ es acíclico. \\
            
            Si el algoritmo rechaza, entonces en algún momento queda un subgrafo no vacío sin ninguna fuente. Por el apartado (a), dicho subgrafo no puede ser acíclico, ya que todo grafo dirigido finito acíclico tiene una fuente. Por tanto, ese subgrafo tiene un ciclo dirigido. Como dicho subgrafo procede de eliminar nodos y arcos de $G$, ese ciclo también estaba en $G$. Luego $G$ no es acíclico.

            Finalmente, el algoritmo tiene tiempo polinómico. En cada iteración se puede buscar una fuente recorriendo los nodos y los arcos del grafo restante. Como en cada iteración se elimina un nodo, hay a lo sumo $n$ iteraciones. Por tanto, el tiempo total está acotado por un polinomio en el tamaño de la entrada.
    \end{enumerate}
\end{ejercicio}

\begin{ejercicio}~
    \begin{enumerate}
        \item[(a)] Demostrar que un grafo es \textit{bipartito} (se puede dividir en dos partes, que no son necesariamente iguales, de manera que 
        todas las aristas van de una parte a otra) si y solo si todos sus ciclos son de longitud par. \\

        Sea $G$ un grafo. Queremos demostrar que, llamando $P$ a la propiedad de que todos los ciclos de $G$ tienen longitud par, entonces
        $$
        G \text{ es bipartito} \Longleftrightarrow \text{se cumple } P.
        $$

        \begin{itemize}
            \item[$\Longrightarrow)$] Supongamos primero que $G=(V,E)$ es bipartito. Entonces existen dos subconjuntos $V_1,V_2 \subset V$ tales que
            $$
            V=V_1\cup V_2,\qquad V_1\cap V_2=\emptyset,
            $$
            y toda arista de $G$ tiene un extremo origen en $V_1$ y el otro destino en $V_2$. Sea
            $$
            v_1\to v_2\to \cdots \to v_k\to v_1
            $$
            un ciclo de $G$, para cierto $k \in \N$. Sin pérdida de generalidad, supongamos que $v_1\in V_1$. Como toda arista va de una parte a la otra, entonces $v_2\in V_2$, $v_3\in V_1$, $v_4\in V_2$, y así sucesivamente. Por tanto, los vértices en posiciones impares pertenecen a $V_1$ y los vértices en posiciones pares pertenecen a $V_2$.
            Como el último arco del ciclo es $(v_k,v_1)$, el vértice $v_k$ debe estar en una parte distinta de la de $v_1$. Como $v_1\in V_1$, 
            necesariamente $v_k\in V_2$. Por lo anterior, esto ocurre si y solo si $k$ es par. Luego todo ciclo de $G$ tiene longitud par.

            \item[$(\Longleftarrow$] Recíprocamente, supongamos que todos los ciclos de $G$ tienen longitud par. Veamos que $G$ es bipartito.
            Podemos suponer primero que $G$ es conexo. Tomamos un vértice cualquiera $s\in V$ y definimos
            $$
            V_1=\{v\in V : \text{existe un camino de longitud par desde } s \text{ hasta } v\},
            $$
            $$
            V_2=\{v\in V : \text{existe un camino de longitud impar desde } s \text{ hasta } v\}.
            $$

            Veamos que esta definición es correcta, es decir, que ningún vértice puede pertenecer a la vez a $V_1$ y a $V_2$. Supongamos, por reducción al absurdo, que existe un vértice $v\in V_1\cap V_2$. Entonces existe un camino de longitud par desde $s$ hasta $v$ y también un camino de longitud impar desde $s$ hasta $v$.
            Tomando ambos caminos, y eliminando las partes comunes si fuera necesario, se obtiene un ciclo de longitud impar. Esto contradice la hipótesis de que todos los ciclos de $G$ tienen longitud par. Por tanto, $V_1\cap V_2=\emptyset$.

            Como $G$ es conexo, todo vértice es alcanzable desde $s$, luego todo vértice pertenece a $V_1$ o a $V_2$. Así,
            $$
            V=V_1\cup V_2.
            $$

            Veamos ahora que toda arista tiene un extremo en $V_1$ y el otro en $V_2$. Sea $(u,v)\in E$. Si $u$ y $v$ pertenecieran al mismo 
            conjunto ($V_1$ o $V_2$), entonces existirían caminos desde $s$ hasta $u$ y desde $s$ hasta $v$ de la misma paridad. 
            Al añadir la arista $(u,v)$, se obtendría un ciclo de longitud impar, lo cual es una contradicción
            con la hipótesis de que todos los ciclos tienen longitud par. Por tanto, toda arista une un vértice de $V_1$ con un vértice de $V_2$.
            Luego $G$ es bipartito. Si $G$ no es conexo, aplicamos el mismo razonamiento a cada componente conexa y unimos las dos partes obtenidas en cada componente. Por tanto, también en el caso general $G$ es bipartito.

        \end{itemize}

        \item[(b)] Utilizar (a) para describir un algoritmo en tiempo polinómico para comprobar si un grafo es bipartito. \\
        
            El algoritmo es el siguiente:

            \begin{algorithm}
                \caption{Decidir si un grafo es bipartito}
                \label{algorithm:bipartito}
                \begin{algorithmic}
                    \State \textbf{Entrada:} Un grafo $G=(V,E)$

                    \State Inicialmente, todos los vértices están sin colorear.

                    \For{cada vértice $s\in V$ sin colorear}
                        \State Colorear $s$ con el color $1$.
                        \State Inicializar una cola con el vértice $s$.

                        \While{la cola no esté vacía}
                            \State Extraer un vértice $u$ de la cola.

                            \For{cada vecino $v$ de $u$}
                                \If{$v$ no está coloreado}
                                    \State Colorear $v$ con el color distinto al de $u$.
                                    \State Añadir $v$ a la cola.
                                \ElsIf{$v$ tiene el mismo color que $u$}
                                    \State Rechazar, es decir, responder que $G$ no es bipartito.
                                \EndIf
                            \EndFor
                        \EndWhile
                    \EndFor

                    \State Aceptar, es decir, responder que $G$ es bipartito.
                \end{algorithmic}
            \end{algorithm}

            Veamos que el algoritmo es correcto. Si el algoritmo acepta, entonces ha asignado a cada vértice uno de dos colores y ha comprobado que ninguna arista une dos vértices del mismo color. Por tanto, tomando como $V_1$ los vértices de color $1$ y como $V_2$ los vértices de color $2$, 
            obtenemos una partición de $V$ tal que todas las aristas van de una parte a la otra. Luego $G$ es bipartito. \\

            Si el algoritmo rechaza, entonces ha encontrado una arista $(u,v) \in E$ cuyos extremos tienen el mismo color. Como los colores se han asignado alternando según la distancia desde el vértice inicial de la componente, que 
            $u$ y $v$ tengan el mismo color significa que hay caminos desde el vértice inicial hasta $u$ y hasta $v$ de la misma paridad. 
            Al añadir la arista $(u,v)$ se obtiene un ciclo de longitud impar (esto se ha visto en el razonamiento del recíproco en (b)). 
            Por el apartado (a), $G$ no es bipartito.
    \end{enumerate}
\end{ejercicio}

\begin{ejercicio}
    Demostrar que \Pol \text{ es cerrada para la unión y la intersección}. \\

    Sean $L_1,L_2 \in \Pol$. Por definición de $\Pol$, existen dos máquinas de Turing deterministas
    $M_1$ y $M_2$ que deciden $L_1$ y $L_2$, respectivamente, en tiempo polinómico. Es decir, existen polinomios $p_1$ y $p_2$ tales que, para toda palabra $x$,
    $M_1$ decide si $x \in L_1$ en tiempo a lo sumo $p_1(|x|)$, y $M_2$ decide si $x \in L_2$ en tiempo a lo sumo $p_2(|x|)$. Veamos primero que 
    $L_1 \cup L_2 \in \Pol$. Construimos una máquina determinista $M_{\cup}$:

    \begin{algorithm}
    \caption{MT $M_{\cup}$}
    \begin{algorithmic}
        \State \textbf{Entrada:} Una palabra $x$

        \State Ejecuta $M_1$ sobre $x$.
        \If{$M_1$ acepta}
            \State $M_{\cup}$ acepta.
        \EndIf

        \State Ejecuta $M_2$ sobre $x$.
        \If{$M_2$ acepta}
            \State $M_{\cup}$ acepta.
        \Else
            \State $M_{\cup}$ rechaza.
        \EndIf
    \end{algorithmic}
    \end{algorithm}

    Entonces $M_{\cup}$ acepta $x$ si y solo si $M_1$ acepta $x$ o $M_2$ acepta $x$, es decir,
    si y solo si $x \in L_1$ o $x \in L_2$. Por tanto, $L(M_{\cup}) = L_1 \cup L_2.$

    Además, su tiempo de ejecución está acotado por
    $$
    p_1(|x|) + p_2(|x|),
    $$
    que es un polinomio. Luego $L_1 \cup L_2 \in \Pol$. Veamos ahora que $L_1 \cap L_2 \in \Pol$. Construimos una máquina determinista $M_{\cap}$
    (\ref{algorithm:cod3}). Entonces $M_{\cap}$ acepta $x$ si y solo si $M_1$ acepta $x$ y $M_2$ acepta $x$, es decir,
    si y solo si $x \in L_1$ y $x \in L_2$. Por tanto,
    $$
    L(M_{\cap}) = L_1 \cap L_2.
    $$

    \begin{algorithm}
    \caption{MT $M_{\cap}$}
    \begin{algorithmic}
        \State \textbf{Entrada:} Una palabra $x$

        \State Ejecuta $M_1$ sobre $x$.
        \If{$M_1$ rechaza}
            \State $M_{\cap}$ rechaza.
        \EndIf

        \State Ejecuta $M_2$ sobre $x$.
        \If{$M_2$ acepta}
            \State $M_{\cap}$ acepta.
        \Else
            \State $M_{\cap}$ rechaza.
        \EndIf
    \end{algorithmic}
    \label{algorithm:cod3}
    \end{algorithm}



    Además, su tiempo de ejecución está acotado por
    $$
    p_1(|x|) + p_2(|x|),
    $$
    que también es un polinomio. Luego $L_1 \cap L_2 \in \Pol$.
\end{ejercicio}



\begin{ejercicio}
    Sea $C$ una clase de funciones de los enteros no negativos en los enteros no negativos. Se dice que $C$ es cerrada para la composición polinómica por la izquierda si $f(n) \in C$ implica que $p(f(n)) = O(g(n))$ para alguna función $g(n) \in C$, donde $p(n)$ es un polinomio cualquiera.
    Se dice que $C$ es cerrada para la composición polinómica por la derecha si $f(n) \in C$ implica que $f(p(n)) = O(g(n))$ para alguna función $g(n) \in C$, donde $p(n)$ es un polinomio cualquiera. \\

    \noindent Determinar cuáles de las siguientes clases de funciones son cerradas para la composición polinómica por la derecha y cuáles lo son por la izquierda. \\
    
    \noindent Nota: en los dos últimos casos se entiende que las funciones se redondean al 
    entero más cercano. En vez de considerar todos los polinomios, como estamos 
    hablando de órdenes de complejidad, será suficiente considerar los polinomios
    de la forma $p(n) = n^l$, donde $l$ es un número natural.

    \begin{enumerate}
        \item[(a)] $C=\{n^k:k>0\}$. \\

            Sea $f(n)=n^k \in C$ y sea $p(n)=n^l$. Por la izquierda,
            $$
            p(f(n))=(n^k)^l=n^{kl}.
            $$
            Como $kl > 0$, tomando $g(n)=n^{kl}\in C$, se tiene
            $$
            p(f(n))=O(g(n)).
            $$
            Luego $C$ es cerrada por la izquierda. Por la derecha,
            $$
            f(p(n))=f(n^l)=(n^l)^k=n^{lk}.
            $$
            Como $lk>0$, tomando $g(n)=n^{lk}\in C$, se tiene
            $$
            f(p(n))=O(g(n)).
            $$
            Luego $C$ es cerrada por la derecha. Por tanto, $C$ es cerrada por la izquierda y por la derecha.

        \item[(b)] $C=\{k n:k>0\}$. \\

            Sea $f(n)=kn\in C$. Tomamos $p(n)=n^2$. Entonces, por la izquierda,
            $$
            p(f(n))=(kn)^2=k^2n^2.
            $$
            Si existiera $g(n)=an\in C$ tal que $k^2n^2=O(an)$, entonces 
            tendríamos $n^2=O(n)$, lo cual es falso. Por tanto, $C$ no es 
            cerrada por la izquierda. Por la derecha,
            $$
            f(p(n))=f(n^2)=kn^2.
            $$
            Si existiera $g(n)=an\in C$ tal que $kn^2=O(an)$, entonces tendríamos $n^2=O(n)$, lo cual es falso. Por tanto, $C$ no es cerrada por la derecha.
            Por tanto, $C$ no es cerrada ni por la izquierda ni por la derecha.

        \item[(c)] $C=\{k^n:k>0\}$. \\

            Sea $f(n)=k^n\in C$ y sea $p(n)=n^l$. Por la izquierda,
            $$
            p(f(n))=(k^n)^l=k^{ln}=(k^l)^n.
            $$
            Como $k^l>0$, tomando $g(n)=(k^l)^n\in C$, se tiene
            $$
            p(f(n))=O(g(n)).
            $$
            Luego $C$ es cerrada por la izquierda. Veamos que no es cerrada por 
            la derecha. Tomamos $f(n)=2^n\in C$ y $p(n)=n^2$. Entonces
            $$
            f(p(n))=f(n^2)=2^{n^2}.
            $$
            Si $C$ fuera cerrada por la derecha, existiría $g(n)=a^n\in C$ tal que
            $$
            2^{n^2}=O(a^n).
            $$
            Pero
            $$
            \frac{2^{n^2}}{a^n}=\left(\frac{2^n}{a}\right)^n \longrightarrow \infty,
            $$
            para cualquier constante $a>0$. Por tanto, $2^{n^2}$ no es $O(a^n)$ para ningún $a>0$.
            Luego $C$ es cerrada por la izquierda, pero no por la derecha.

        \item[(d)] $C=\{2^{n^k}:k>0\}$. \\

            Sea $f(n)=2^{n^k}\in C$ y sea $p(n)=n^l$. Por la izquierda,
            $$
            p(f(n))=(2^{n^k})^l=2^{l n^k}.
            $$
            Como
            $$
            \lim_{n \to \infty} \frac{l n^k}{n^{k+1}}
            =
            \lim_{n \to \infty} \frac{l}{n}
            =
            0,
            $$
            se tiene
            $$
            2^{l n^k}=O(2^{n^{k+1}}).
            $$
            Tomando $g(n)=2^{n^{k+1}}\in C$, se deduce que $C$ es cerrada por la izquierda.
            Por la derecha,
            $$
            f(p(n))=f(n^l)=2^{(n^l)^k}=2^{n^{lk}}.
            $$
            Como $lk>0$, tomando $g(n)=2^{n^{lk}}\in C$, se tiene
            $$
            f(p(n))=O(g(n)).
            $$
            Luego $C$ es cerrada por la derecha.
            Por tanto, $C$ es cerrada por la izquierda y por la derecha.

        \item[(e)] $C=\{\log^k(n):k>0\}$. \\

            Sea $f(n)=\log^k(n)\in C$ y sea $p(n)=n^l$. Por la izquierda,
            $$
            p(f(n))=(\log^k(n))^l=\log^{kl}(n).
            $$
            Como $kl>0$, tomando $g(n)=\log^{kl}(n)\in C$, se tiene
            $$
            p(f(n))=O(g(n)).
            $$
            Luego $C$ es cerrada por la izquierda. Por la derecha,
            $$
            f(p(n))=f(n^l)=\log^k(n^l)=(l\log(n))^k=l^k\log^k(n).
            $$
            Por tanto,
            $$
            f(p(n))=O(\log^k(n)).
            $$
            Tomando $g(n)=\log^k(n)\in C$, se deduce que $C$ es cerrada por la derecha.
            Por tanto, $C$ es cerrada por la izquierda y por la derecha.

        \item[(f)] $C=\{\log(n)\}$. \\

            Sea $f(n)=\log(n)$. Tomamos $p(n)=n^2$. Por la izquierda,
            $$
            p(f(n))=(\log(n))^2.
            $$
            Si $C$ fuera cerrada por la izquierda, tendría que cumplirse
            $$
            (\log(n))^2=O(\log(n)),
            $$
            lo cual es falso. Por tanto, $C$ no es cerrada por la izquierda.
            Por la derecha, sea $p(n)=n^l$. Entonces
            $$
            f(p(n))=f(n^l)=\log(n^l)=l\log(n).
            $$
            Así,
            $$
            f(p(n))=O(\log(n)).
            $$
            Como $\log(n)\in C$, se deduce que $C$ es cerrada por la derecha.
            Por tanto, $C$ es cerrada por la derecha, pero no por la izquierda.
    \end{enumerate}
\end{ejercicio}



\begin{ejercicio}
    Sea $L = \{wcw : w \in \{0,1\}^*\}.$ Demostrar que $L$ se puede reconocer en espacio $\log(n)$. \\

    Sea $x$ una palabra de entrada y sea $n = |x|$. Construimos una MT determinista con una cinta de entrada de solo lectura y una cinta de trabajo.
    La máquina funciona de la siguiente forma:

    \begin{algorithm}
        \caption{MT $M$}
        \label{algorithm:wcw-log}
        \begin{algorithmic}
            \State \textbf{Entrada:} Una palabra $x \in \{0,1,c\}^{*}$

            \State Calcula $n = |x|$ usando un contador binario.
            \State Recorre $x$ y comprueba que aparece exactamente un símbolo $c$.
            \If{no aparece exactamente un símbolo $c$}
                \State $M$ rechaza.
            \EndIf

            \State Sea $k$ la posición en la que aparece $c$.
            \If{$n \neq 2k-1$}
                \State $M$ rechaza.
            \EndIf

            \For{$i = 1,\dots,k-1$}
                \State Lee el símbolo de $x$ en la posición $i$ y lo guarda.
                \State Lee el símbolo de $x$ en la posición $k+i$ y lo compara con el anterior.
                \If{los símbolos son distintos}
                    \State $M$ rechaza.
                \EndIf
            \EndFor

            \State $M$ acepta.
        \end{algorithmic}
    \end{algorithm}

    Veamos que el algoritmo es correcto. Si $M$ acepta, entonces $x$ tiene exactamente un símbolo $c$, situado en una posición $k$ tal que $n = 2k-1$. 
    Por tanto, a la izquierda y a la derecha de $c$ hay exactamente $k-1$ símbolos. Además, para cada $i = 1,\dots,k-1$, el símbolo en la posición $i$ 
    coincide con el símbolo en la posición $k+i$. Luego ambas mitades son iguales, y por tanto $x = wcw$ para cierta palabra $w \in \{0,1\}^{*}$. \\

    Recíprocamente, si $x \in L$, entonces $x = wcw$ para alguna palabra $w \in \{0,1\}^{*}$. En ese caso, $x$ tiene exactamente un símbolo $c$, 
    las dos partes a izquierda y derecha de $c$ tienen la misma longitud, y todos los símbolos correspondientes coinciden. Por tanto, el algoritmo no rechaza en 
    ningún paso y acepta. \\

    Finalmente, analizamos el espacio. La máquina solo necesita guardar un número constante de contadores, como $n$, $k$, $i$ y el contador auxiliar usado para 
    moverse hasta una posición concreta de la entrada. Todos esos contadores están acotados por $n$, luego cada uno ocupa $O(\log n)$ casillas. Además, solo se guarda 
    un número constante de símbolos de $\{0,1,c\}$. Por tanto, el espacio total usado en la cinta de trabajo es $O(\log n)$. Concluimos que
    $$L \in \ESP{\log n}.$$
\end{ejercicio}



\begin{ejercicio}
    Demostrar que si $L$ está en \Pol, entonces $L^*$ también está en \Pol. \\

    Como $L \in \Pol$, existe un algoritmo determinista $A$ que decide $L$ en tiempo
    polinómico. Sea $p(n)$ un polinomio tal que $A$ decide si una palabra de longitud
    $n$ pertenece a $L$ en tiempo $O(p(n))$. Construimos un algoritmo 
    determinista para decidir $L^*$. Sea $x = x_1 \cdots x_n$
    la palabra de entrada. Usaremos programación dinámica. Definimos, para cada
    $i = 0,\dots,n$,
    $$D[i] = 1 \iff x_1 \cdots x_i \in L^*.$$
    Inicializamos $D[0] = 1$, ya que $\varepsilon \in L^*$. Para calcular 
    $D[i]$, con $i \geq 1$, usamos la equivalencia
    $$x_1 \cdots x_i \in L^*
    \iff \exists j \in \{0,\dots,i-1\} \text{ tal que } D[j] = 1
    \text{ y } x_{j+1} \cdots x_i \in L.$$

    Por tanto, el algoritmo es:

    \begin{algorithm}
        \caption{Decisor para $L^*$}
        \begin{algorithmic}
            \State \textbf{Entrada:} Una palabra $x = x_1 \cdots x_n$

            \State $D[0] \gets 1$

            \For{$i = 1,\dots,n$}
                \State $D[i] \gets 0$
                \For{$j = 0,\dots,i-1$}
                    \If{$D[j] = 1$ y $A$ acepta $x_{j+1} \cdots x_i$}
                        \State $D[i] \gets 1$
                    \EndIf
                \EndFor
            \EndFor

            \If{$D[n] = 1$}
                \State acepta
            \Else
                \State rechaza
            \EndIf
        \end{algorithmic}
    \end{algorithm}

    Veamos que el algoritmo es correcto. Por definición, $D[0] = 1$, pues
    $\varepsilon \in L^*$. Supongamos calculados correctamente $D[0],\dots,D[i-1]$.
    Entonces $D[i] = 1$ si y solo si existe una posición $j < i$ tal que el prefijo
    $x_1 \cdots x_j$ pertenece a $L^*$ y el bloque final $x_{j+1}\cdots x_i$ pertenece
    a $L$. Esto es exactamente decir que $x_1 \cdots x_i$ es concatenación de palabras
    de $L$, es decir, que $x_1 \cdots x_i \in L^*$. Por inducción, $D[n] = 1$ si y
    solo si $x \in L^*$. \\

    Finalmente, analizamos el tiempo. Hay $O(n^2)$ pares $(i,j)$, y para cada uno se
    ejecuta el algoritmo $A$ sobre una palabra de longitud menor o igual que $n$, lo que
    computacionalmente es de orden $O(p(n))$. Por tanto, el tiempo total está acotado por
    $$O(n^2p(n)),$$ que es polinómico. Luego $L^* \in \Pol$.
\end{ejercicio}



% TODO: revisar el siguiente

\begin{ejercicio}
    Demostrar que si $L$ está en \NP, entonces $L^*$ también está en \NP. \\

    Como $L \in \NP$, existe un verificador determinista $V$ en tiempo polinómico y
    un polinomio $p$ tales que, para toda palabra $u$,
    $$u \in L \iff \exists c,\ |c| \leq p(|u|), \text{ tal que } V(u,c) \text{ acepta}.$$

    Veamos que $L^* \in \NP$. Sea $x = x_1 \cdots x_n$ una palabra de entrada.
    Un certificado de que $x \in L^*$ consiste en una descomposición de $x$ en bloques
    no vacíos
    $$x = u_1u_2\cdots u_k,$$
    junto con un certificado $c_i$ para cada bloque $u_i$, demostrando que
    $u_i \in L$.

    Equivalentemente, el certificado contiene posiciones
    $$0 = i_0 < i_1 < \cdots < i_k = n,$$
    de forma que
    $$u_j = x_{i_{j-1}+1}\cdots x_{i_j},$$
    y certificados $c_1,\dots,c_k$ tales que $V(u_j,c_j)$ acepta para todo
    $j = 1,\dots,k$.

    El verificador para $L^*$ funciona así:

    \begin{algorithm}
        \caption{Verificador para $L^*$}
        \begin{algorithmic}
            \State \textbf{Entrada:} Una palabra $x = x_1 \cdots x_n$ y un certificado
            formado por $0 = i_0 < i_1 < \cdots < i_k = n$ y palabras $c_1,\dots,c_k$

            \If{$x = \varepsilon$}
                \State acepta
            \EndIf

            \For{$j = 1,\dots,k$}
                \State Sea $u_j = x_{i_{j-1}+1}\cdots x_{i_j}$
                \State Ejecuta $V(u_j,c_j)$
                \If{$V(u_j,c_j)$ rechaza}
                    \State rechaza
                \EndIf
            \EndFor

            \State acepta
        \end{algorithmic}
    \end{algorithm}

    Veamos que es correcto. Si $x \in L^*$, entonces existe una descomposición
    $$x = u_1u_2\cdots u_k$$
    con $u_j \in L$ para todo $j$. Como $L \in \NP$, para cada $u_j$ existe un
    certificado $c_j$ que hace aceptar a $V(u_j,c_j)$. Por tanto, el verificador
    anterior acepta $x$.

    Recíprocamente, si el verificador acepta, entonces cada bloque $u_j$ pertenece a
    $L$, ya que $V(u_j,c_j)$ acepta. Por tanto,
    $$x = u_1u_2\cdots u_k \in L^*.$$

    Finalmente, veamos que el tamaño del certificado y el tiempo de verificación son
    polinómicos. Como los bloques son no vacíos, se tiene $k \leq n$. Además,
    $|u_j| \leq n$, luego cada certificado $c_j$ tiene tamaño como mucho $p(n)$.
    Por tanto, el tamaño total del certificado está acotado por
    $$O(np(n)),$$
    que es polinómico. Además, el verificador ejecuta a lo sumo $n$ veces un algoritmo
    polinómico sobre palabras de longitud menor o igual que $n$. Luego el tiempo total
    también es polinómico.

    Así, $L^* \in \NP$.
\end{ejercicio}



\begin{ejercicio}
    Demostrar que \NP \text{ es cerrada para la unión y la intersección}. \\

    Sean $L_1,L_2 \in \NP$. Por definición de $\NP$, existen dos máquinas de Turing no deterministas
    $M_1$ y $M_2$ que deciden $L_1$ y $L_2$, respectivamente, en tiempo polinómico.

    Es decir, existen polinomios $p_1$ y $p_2$ tales que, para toda palabra $x$,
    $M_1$ decide si $x \in L_1$ en tiempo a lo sumo $p_1(|x|)$, y $M_2$ decide si
    $x \in L_2$ en tiempo a lo sumo $p_2(|x|)$.
    Veamos primero que $L_1 \cup L_2 \in \NP$. Construimos una máquina de Turing no determinista
    $M_{\cup}$:

    \begin{algorithm}
    \caption{MTND $M_{\cup}$}
    \begin{algorithmic}
        \State \textbf{Entrada:} Una palabra $x$

        \State Elegimos de forma no determinista $i \in \{1,2\}$.

        \If{$i=1$}
            \State Simulamos $M_1$ sobre $x$.
        \Else
            \State Simulamos $M_2$ sobre $x$.
        \EndIf

        \State $M_{\cup}$ acepta si y solo si la máquina simulada acepta.
    \end{algorithmic}
    \end{algorithm}

    Entonces $M_{\cup}$ acepta $x$ si y solo si existe una rama de cómputo que acepta en $M_1$
    o existe una rama de cómputo que acepta en $M_2$. Por tanto,
    $$
    x \in L(M_{\cup}) \iff x \in L_1 \cup L_2.
    $$

    Además, el tiempo de ejecución de $M_{\cup}$ está acotado por
    $$
    \max\{p_1(|x|),p_2(|x|)\} + O(1),
    $$
    que es un polinomio. Luego $L_1 \cup L_2 \in \NP$. \\

    Veamos ahora que $L_1 \cap L_2 \in \NP$. Construimos una máquina de Turing no determinista
    $M_{\cap}$:

    \begin{algorithm}
    \caption{MTND $M_{\cap}$}
    \begin{algorithmic}
        \State \textbf{Entrada:} Una palabra $x$

        \State Simulamos $M_1$ sobre $x$.
        \If{$M_1$ rechaza}
            \State $M_{\cap}$ rechaza.
        \EndIf

        \State Simulamos $M_2$ sobre $x$.
        \If{$M_2$ acepta}
            \State $M_{\cap}$ acepta.
        \Else
            \State $M_{\cap}$ rechaza.
        \EndIf
    \end{algorithmic}
    \end{algorithm}

    Esta máquina acepta $x$ si y solo si existe una rama aceptante de $M_1$ sobre $x$ y,
    después, existe una rama aceptante de $M_2$ sobre $x$. Por tanto,
    $$
    x \in L(M_{\cap}) \iff x \in L_1 \cap L_2.
    $$

    Además, el tiempo de ejecución está acotado por
    $$
    p_1(|x|)+p_2(|x|),
    $$
    que es un polinomio. Luego $L_1 \cap L_2 \in \NP$.
\end{ejercicio}

\begin{ejercicio}
    Demostrar que si $\NP \neq \CoNP$, entonces $\Pol \neq \NP$. \\

    Demostraremos el contrarrecíproco, es decir, si $\Pol = \NP$, entonces
    $\NP = \CoNP$. Lo probamos por doble contenido.
    
    \begin{itemize}
        \item[$\subseteq)$] Sea $L \in \NP$. Como $\NP = \Pol$, se tiene que $L \in \Pol$, y como $\Pol$ es cerrada
            por complementarios, $\overline{L} \in \Pol$. De nuevo, como $\Pol = \NP$, se deduce que
            $\overline{L} \in \NP$. Por definición de $\CoNP$, esto implica que $L \in \CoNP$. Por tanto,
            $$\NP \subseteq \CoNP.$$
        \item[$(\supseteq$] Recíprocamente, sea $L \in \CoNP$. Entonces, por definición de $\CoNP$, se tiene que
            $\overline{L} \in \NP$. Como $\NP = \Pol$, resulta que $\overline{L} \in \Pol$. Usando otra vez
            que $\Pol$ es cerrada por complementario, obtenemos que $L \in \Pol$. Finalmente, como
            $\Pol = \NP$, concluimos que $L \in \NP$. Por tanto,
            $$\CoNP \subseteq \NP.$$
    \end{itemize}
\end{ejercicio}

\begin{ejercicio}
    Supongamos que una entrada es una palabra de paréntesis. 
    Demostrar que determinar si están emparejados y anidados correctamente 
    está en \Loga. Lo están $(()())$ y $((()))$, pero no lo están $)()($ ni 
    $(()))($. \\

    Sea $L=\{w\in \{(,)\}^{*} : w \text{ está correctamente emparejada y anidada}\}.$
    Vamos a demostrar que $L\in \Loga$ construyendo un algoritmo determinista que decide $L$ usando espacio logarítmico. \\

    Dada una entrada $w\in \{(,)\}^{*}$, el Algoritmo (\ref{algorithm:cod10}) recorre $w$ de izquierda a derecha manteniendo un contador $c$, 
    que representa el número de paréntesis abiertos que todavía no han sido cerrados. \\

    \begin{algorithm}
        \caption{Algoritmo para paréntesis bien anidados}
        \label{algorithm:cod10}
        \begin{algorithmic}
            \State \textbf{Entrada:} Una palabra $w\in \{(,)\}^{*}$
            \State $c \gets 0$

            \For{cada símbolo $a$ de $w$, de izquierda a derecha}
                \If{$a = ($}
                    \State $c \gets c+1$
                \EndIf

                \If{$a = )$}
                    \If{$c=0$}
                        \State Rechazar
                    \Else
                        \State $c \gets c-1$
                    \EndIf
                \EndIf
            \EndFor

            \If{$c=0$}
                \State Aceptar
            \Else
                \State Rechazar
            \EndIf
        \end{algorithmic}
    \end{algorithm}

    Veamos que el algoritmo es correcto. Si el algoritmo acepta, entonces nunca ha leído un paréntesis cerrado sin 
    que hubiese antes un paréntesis abierto disponible, pues en tal caso se tendría $c=0$ al leer un símbolo $)$ y el algoritmo rechazaría. 
    Además, al final se tiene $c=0$, luego todos los paréntesis abiertos han sido cerrados. Por tanto, $w$ está correctamente emparejada y anidada. \\

    Recíprocamente, si $w$ está correctamente emparejada y anidada, entonces en ningún prefijo de $w$ hay más 
    paréntesis cerrados que abiertos. Por tanto, el contador nunca intenta bajar de $0$. Además, al final hay tantos paréntesis abiertos como 
    cerrados, luego $c=0$. Así, el algoritmo acepta. \\

    Finalmente, si $n=|w|$, el contador $c$ nunca toma valores mayores que $n$. Por tanto, puede almacenarse usando $\lceil \log(n+1)\rceil$
    bits. Aparte de este contador, el algoritmo solo usa espacio constante. Luego el espacio total usado es $O(\log n)$.
    Por tanto, $L$ se decide en espacio logarítmico, es decir, $L\in \Loga$.
\end{ejercicio}



\begin{ejercicio}
    Consideremos el problema anterior, pero ahora con dos tipos de 
    paréntesis: $(,)$ y $[,]$. Demostrar que determinar si están bien 
    escritos también está en \Loga. Aquí $([]())$ está permitido, pero 
    $([)]$ no. Demostrar que si la entrada se lee de izquierda a derecha 
    sin volver hacia atrás, requieren $O(n)$ de memoria. \\

    Sea $L=\{w\in\{(,),[,]\}^{*} : w \text{ está bien escrita}\}.$
    Primero comprobamos, como en el ejercicio anterior, que $w$ está correctamente emparejada ignorando el tipo de paréntesis. 
    Es decir, recorremos la entrada con un contador $c$: sumamos $1$ al leer $($ o $[$, restamos $1$ al leer $)$ o $]$, rechazamos si 
    $c$ intenta hacerse negativo, y al final exigimos $c=0$. Esto usa espacio $O(\log n)$. \\

    Falta comprobar que cada apertura se cierra con el símbolo del mismo tipo. Para ello, para cada posición $i$ de la entrada que contenga un símbolo de apertura, 
    buscamos su cierre correspondiente. Recorremos la entrada desde la posición $i+1$ hacia la derecha con un contador $c$ inicializado a $1$. Cada vez que leemos una 
    apertura aumentamos $c$, y cada vez que leemos un cierre disminuimos $c$. El primer símbolo que hace que $c=0$ es el cierre correspondiente a la apertura situada en $i$. \\
    
    Así, si en la posición $i$ hay $($, comprobamos que dicho cierre sea $)$; si en la posición $i$ hay $[$, comprobamos que dicho cierre sea $]$. 
    Si alguna comprobación falla, rechazamos. Si todas son correctas, aceptamos. \\

    El algoritmo es correcto porque, una vez comprobado que la palabra está bien emparejada ignorando tipos, cada símbolo de apertura tiene un único símbolo de cierre correspondiente. Por tanto, basta comprobar que ambos son del mismo tipo.

    En cuanto al espacio, solo se necesita guardar la posición $i$, el contador $c$ y una cantidad constante de información auxiliar. Como $i\leq n$ y $c\leq n$, ambos se almacenan con $O(\log n)$ bits. Por tanto, el espacio total usado es $O(\log n)$, y se concluye que
    $L\in \Loga.$ \\

    Por último, si la entrada se lee de izquierda a derecha sin volver hacia atrás, se necesita guardar los tipos de paréntesis abiertos todavía no cerrados. Esto se hace con una pila. En el peor caso, la entrada puede comenzar con $n/2$ símbolos de apertura, por ejemplo una palabra de la forma
    $$
    x_1x_2\cdots x_{n/2},
    $$
    donde cada $x_i$ es $($ o $[$. 
    Para poder comprobar después que los cierres aparecen en el orden y con el tipo correcto, hay que recordar esos $n/2$ tipos. 
    Luego se requiere memoria $O(n)$.
\end{ejercicio}

% TODO: revisar esto y ver cómo redactarlo mejor

\begin{ejercicio}
    Demostrar que el problema del palíndromo requiere $O(n)$ de memoria si no se puede volver hacia atrás. \\

    En realidad, lo que se demuestra es que requiere memoria lineal, es decir, $\Omega(n)$, y como existe un algoritmo que usa $O(n)$ memoria, el problema requiere $\Theta(n)$ memoria.

    Consideremos palabras sobre el alfabeto $\{0,1\}$ y supongamos que existe un algoritmo que lee la entrada una sola vez, de izquierda a derecha, sin volver hacia atrás, y que decide si una palabra es palíndroma usando $s(n)$ bits de memoria.

    Tomemos palabras de longitud $n=2m$. Después de leer los primeros $m$ símbolos, el algoritmo ha leído una palabra $u \in \{0,1\}^m$. Hay $2^m$ posibles palabras $u$.

    Si el algoritmo usara menos de $m$ bits de memoria, entonces tendría menos de $2^m$ configuraciones posibles de memoria. Por el principio del palomar, existirían dos palabras distintas $u,v \in \{0,1\}^m$ tales que, después de leer $u$ y después de leer $v$, el algoritmo queda exactamente en la misma configuración de memoria.

    Como el algoritmo no puede volver hacia atrás, a partir de ese momento se comportará igual ante cualquier continuación. En particular, consideremos la continuación $u^R$, donde $u^R$ es la palabra $u$ escrita al revés.

    Entonces:
    $$
    uu^R
    $$
    es un palíndromo, pero
    $$
    vu^R
    $$
    no es un palíndromo, ya que $v \neq u$.

    Sin embargo, el algoritmo llega a la misma configuración tras leer $u$ y $v$, y después recibe la misma continuación $u^R$. Por tanto, daría la misma respuesta para $uu^R$ y para $vu^R$, contradicción.

    Luego el algoritmo necesita distinguir las $2^m$ posibles primeras mitades, por lo que necesita al menos $m$ bits de memoria. Como $n=2m$, esto implica que necesita
    $$
    \Omega(n)
    $$
    memoria.

    Por otro lado, existe un algoritmo que guarda toda la palabra, o al menos su primera mitad, y luego compara con la segunda mitad. Este algoritmo usa $O(n)$ memoria.

    Por tanto, si la entrada se lee de izquierda a derecha sin volver hacia atrás, el problema del palíndromo requiere memoria lineal:
    $$
    \Theta(n).
    $$
\end{ejercicio}



\begin{ejercicio}
    Demostrar que $\NP \neq \ESP{n}$. Nota: Este ejercicio es muy complicado. Como sugerencia, para probar que 
    las clases son diferentes, probar que $\NP$ cumple la propiedad de ser cerrada para reducciones en $L$ 
    (si $P_1 \propto P_2$ y $P_2$ está en $\NP$, entonces $P_1$ está en $\NP$), mientras que $\ESP{n}$, no. \\

    Sea $P_{*}$ la propiedad ``ser una clase de complejidad cerrada por reducciones en $\Loga$''. 
    Recordemos que $\Loga = \ESP{\log(n)}$ (clase de espacio logarítmico determinista). 
    La demostración consiste en probar dos cosas:
    \begin{enumerate}
        \item $\NP$ cumple $P_{*}$.
        \item $\ESP{n}$ \textbf{no} cumple $P_{*}$.
    \end{enumerate}
    Consecuentemente, deberá ser $\NP \neq \ESP{n}$, ya que si fuera $\NP = \ESP{n}$, en particular, 
    ambas clases deberían cumplir $P_{*}$, cosa que vamos a concluir que no sucede. \\ 
    
    Consideramos la siguiente definición y resultados antes de demostrar:
    
    \begin{definicion}\label{definicion:reduc}
        Un problema de decisión $P_1(x)$ es reducible a un problema de
        decisión $P_2(y)$ donde $x \in X$ e $y \in Y$, denotado por $P_1(x) \propto P_2(y)$ si y solo si 
        existe un algoritmo determinista que en espacio logarítmico calcula una función $\REDU : X \to Y$ 
        de tal manera que $P_1(x) = P_2 (\REDU(x))$.
    \end{definicion}

    \begin{lema}\label{lema:tpol}
        Sea $\REDU: X \to Y$ una función calculable por una MT
        determinista en espacio $O(\log n)$, con $n = |x|$ es la longitud de la entrada $x \in X$, con
        una cinta de entrada de solo lectura, una cinta de trabajo y una cinta de salida de solo escritura que no 
        vuelve hacia atrás en ningún momento. Entonces $\REDU(x)$ se calcula en tiempo polinómico y existe $d \in \N : |\REDU(x)| \leq n^d$.
    \end{lema} 
    
    \begin{proof}
        Sea $M$ una MT determinista que calcula $\mathrm{REDU}$ usando espacio
        $O(\log n)$. Una configuración de $M$ queda 
        determinada por: $$(q,i,u,v)$$
        donde $q$ es el estado, $i$ es la posición del cabezal en la cinta de
        entrada, y $u,v$ describen el contenido de la cinta de trabajo antes y
        después del cabezal. Como $M$ usa espacio $O(\log n)$, existe una constante
        $c > 0$ tal que la longitud total de $u$ y $v$ es menor o igual que
        $c\log n$. Por tanto, el número de configuraciones posibles está acotado por
        $$|Q|\cdot n\cdot |B|^{c\log n}.$$
        Como $|Q|$ y $|B|$ son constantes de la máquina,
        $$|Q|\cdot n\cdot |B|^{c\log n} \quad \text{tiene orden } \quad O\left(n\cdot |B|^{c\log n}\right).
        $$Ahora usando que, para $a > 0$ y $n>0$ se tiene que $a^{\log(n)} = n^{\log(a)}$, entonces, si $a = |B| > 0$:
        $$|B|^{c\log n} = (|B|^{\log(n)})^{c} = (n^{\log |B|})^c = n^{c\log |B|},$$
        luego
        $$|Q|\cdot n\cdot |B|^{c\log n} \quad \text{tiene orden } \quad O(n \cdot n^{c \log |B|}) = O(n^d).
        $$
        para alguna constante $d \in \N$. Es decir, $M$ tiene solo un número polinómico
        de configuraciones posibles. Como $M$ es determinista y siempre termina (por ser REDU calculable total), no puede repetir una
        configuración. Si repitiese una configuración, entraría en un ciclo y no
        terminaría. Por tanto, el número de pasos de cálculo de $M$ está acotado por
        el número de configuraciones posibles, que es polinómico. Finalmente, en cada 
        paso $M$ puede escribir como mucho un símbolo en la cinta
        de salida. Luego, si el número de pasos es polinómico, la longitud de la salida
        también es polinómica. Por tanto,
        $$|\mathrm{REDU}(x)|\leq n^d$$
        y $\mathrm{REDU}$ se calcula en tiempo polinómico (ya que la calcula la MT $M$, que solo da un número
        de pasos polinómico).
    \end{proof}

    \begin{teo}[Teorema de Jerarquía en Espacio]\label{teo:jerarqesp}
        Si $f(n)$ es una función de complejidad propia, entonces 
        $$\ESP{f(n)} \subsetneq \ESP{f(n) \log f(n)}$$ 
    \end{teo}

    \begin{proof}~
        \begin{enumerate}
            \item Para ver que $\NP$ cumple $P_{*}$, hay que ver que, dados $P_1$ y $P_2$ dos problemas, tal que 
            $P_2 \in \NP$ y $P_1 \propto P_2$, entonces $P_1 \in \NP$. Como $P_1 \propto P_2$, por la Definición
            \ref{definicion:reduc} sabemos que existe una reducción $\REDU : X \to Y$ calculable por un algoritmo
            determinista en espacio logarítmico, tal que para toda entrada $x \in X$, entonces
            $P_1(x) = P_2(\REDU(x))$. Como $P_2 \in \NP$, existe un algoritmo no determinista $\ALG_2$ que 
            decide $P_2$ en tiempo polinómico. Es decir, existe un polinomio $p$ tal que, para cualquier entrada
            $y \in Y$, $\ALG_2(y)$ termina en, a lo sumo, tiempo $p(|y|)$. Podemos construir entonces el siguiente
            algoritmo no determinista (ya que delega en $\ALG_2$, que no es determinista) $\ALG_1$ para 
            decidir $P_1$:
            $$
            \begin{array}{l}
            \ALG_1(x):\\
            \quad y := \REDU(x)\\
            \quad \text{return } \ALG_2(y)
            \end{array}
            $$
            Debemos ver que $\ALG_1$ responde correctamente en todos los casos, y que $\ALG_1$ tarda tiempo polinómico:
            \begin{enumerate}
                \item Lo primero es trivial, ya que $P_1(x) = P_2(\REDU(x)) = P_2(y)$. Por tanto, $\ALG_1$ acepta
                $x$ si y solo si $\ALG_2$ acepta $y$, y eso ocurre si y solo si $P_2(y) = \text{SÍ}$, lo cual 
                a su vez equivale a que $P_1(x) = \text{SÍ}$. Lo mismo pasa con NO y no aceptar.
                \item Para ver que $\ALG_1$ tarda tiempo polinómico, como REDU es calculable en espacio
                logarítmico, por el Lema \ref{lema:tpol}, REDU$(x)$ se calcula en tiempo polinómico, y además
                existe $d \in \N$ tal que $|\REDU(x)| \stackrel{(*)}{\leq} n^d$, con $n = |x|$. Por definición de $\ALG_1$, 
                $y = \REDU(x)$, de donde $\ALG_2(y)$ tarda a lo sumo $p(|y|) = p(|\REDU(x)|)$. Como 
                $p$ es un polinomio, $$\exists m \in \N, \quad C_t > 0 : t \in \N 
                \Longrightarrow p(t) \leq C_t t^m$$ 
                luego tenemos que $$p(|\REDU(x)|) \leq C |\REDU(x)|^m \stackrel{(*)}{\leq} C(n^d)^m = C n^{dm}$$
                y $Cn^{dm}$ también es un polinomio en la longitud de la entrada $n = |x|$. \\

                Por lo tanto, el tiempo total de $\ALG_1$ es la suma del tiempo en calcular REDU$(x)$, que por el Lema
                \ref{lema:tpol} es de orden polinómico, y el tiempo para ejecutar 
                $\ALG_2(\REDU(x))$, que, como acabamos de ver, es también de orden polinómico.
            \end{enumerate}
            
            Tenemos entonces que $\ALG_1$ es un algoritmo no determinista que decide correctamente
            $P_1$ en tiempo polinómico. Luego $P_1 \in \NP$. De la arbitrariedad de $P_1$ y $P_2$ 
            con $P_2 \in \NP$ y $P_1 \propto P_2$, se deduce que \NP \text{ cumple} $P_*$.

            \item Para ver que $\ESP{n}$ no cumple $P_{*}$, basta encontrar dos problemas $P_1$ y $P_2$ tales que:
            $P_2 \in \ESP{n}$ y $P_1 \propto P_2$, pero $P_1 \notin \ESP{n}$. \\

            Usaremos el Teorema \ref{teo:jerarqesp}, para la 
            función $f(n) = n$, que es propia (es claramente creciente (de hecho, estrictamente creciente), 
            y además se puede calcular la función $g(u) = f(|u|) = |u|$ en unario usando espacio $O(n)$ (de hecho, $O(\log n) \subseteq O(n)$ usando un contador binario) y tiempo $O(n + n) = O(n)$). Entonces se tiene que 
            $$\ESP{n} \subsetneq \ESP{n \log n}$$
            En tal caso, podemos tomar $P_1 \in \ESP{n \log n} \setminus \ESP{n}$. Definimos $P_2$ como sigue:
            $$P_2 = \{x \# 1^{|x|^2} : x \in P_1\}$$
            Es decir, una entrada de $P_2$ es positiva si tiene la forma $x \# 1^{|x|^2}$ y además 
            $x$ es una entrada positiva de $P_1$. Probamos lo que se quiere en el mismo orden en que se ha enunciado:
            \begin{enumerate}
                \item $P_2 \in \ESP{n}$. Sea $z = x \# 1^{|x|^2}$ una entrada de $P_2$, y sea $N = |z|$. 
                Si $m = |x|$, entonces $N = m + 1 + m^2$. Si $m = 0$, entonces $z = \#$, entonces $N = 1$, y el espacio requerido es constante
                $O(1) \subseteq O(N)$, y en particular estaría en $\ESP{n}$ (aceptamos si $\varepsilon \in P_1$ y rechazamos si $\varepsilon \notin P_1$). Podemos suponer entonces 
                que $m > 0$. Como $P_1 \in \ESP{n \log n}$, decidir si $x \in P_1$
                requiere espacio $O(m \log m)$. Pero, para todo $m > 0$,
                $$m \log m \leq m^2 \leq m + 1 + m^2 = N$$
                Por lo tanto, $O(m \log m) \subseteq O(N)$. Así, la simulación de $\ALG'_1$ (que existe 
                porque $P_1 \in \ESP{n \log n}$) que decide si $x \in P_1$ usa espacio lineal
                respecto al tamaño $N$ de entrada de $P_2$. Además, comprobar que la entrada es de la forma 
                $x \# 1^{|x|^2}$ también se puede hacer en espacio $O(N)$ (recorriendo $z$ y usando contadores binarios acotados por $N$, luego que usen espacio
                $O(\log N) \subseteq O(N)$). Por tanto, $P_2 \in \ESP{n}$.
                \item $P_1 \propto P_2$. Por definición de reducción, debemos construir una función 
                $\REDU' : X \to Y$ tal que $x \in P_1 \iff \REDU'(x) \in P_2$, y además debe ser 
                calculable en espacio logarítmico. Para ello, definimos $$\REDU'(x) = x \# 1^{|x|^2}$$
                Para que sea calculable en espacio logarítmico, buscamos que se escriba la salida símbolo a símbolo
                directamente sin guardar toda la salida en la memoria. El algoritmo sería el siguiente:
                $$
                \begin{array}{l}
                \text{Entrada}: x\\
                \quad \text{Copiar $x$ en la cinta de salida} \\
                \quad \text{Escribir \# en la cinta de salida } \\
                \quad \text{Escribir tantos 1's como indique } |x|^2
                \end{array}
                $$
                La primera operación es $O(1)$ en espacio, ya que no hay que almacenar $x$, y se va copiando
                símbolo a símbolo. La segunda también es $O(1)$ en espacio por el mismo motivo para el 
                símbolo \#, y para la tercera, se obtiene $n = |x|$ por medio de un contador binario, que usa espacio $O(\log n)$. Luego se usan 
                dos contadores binarios $i$ y $j$ para escribir, para cada par $(i,j)$,
                con $1 \leq i,j \leq n$, un 1. Cada contador solo necesita almacenar un número entre $1$ y $n$.
                En binario, necesitaremos $\log n$ bits como máximo por contador. Como son dos contadores, el orden del espacio
                será $O(\log n) + O(\log n) = O(\log n)$. Así, REDU$'$ es calculable en espacio logarítmico.
                Por definición, es claro que $\REDU'$ es una reducción de $P_1$ a $P_2$. Consecuentemente,
                $P_1 \propto P_2$.  
            \end{enumerate}

            Como por definición $P_1 \notin \ESP{n}$, entonces $\ESP{n}$ no cumple $P_{*}$.
        \end{enumerate}
        Tenemos entonces que $\NP$ cumple $P_{*}$ pero $\ESP{n}$ no la cumple, luego, por lo que se ha dicho
        al principio del ejercicio, se tiene que $\NP \neq \ESP{n}$, como se pedía demostrar.
    \end{proof}
\end{ejercicio}
